\conclusion

В ходе выполнения выпускной квалификационной работы была рассмотрена проблема доступа пользователей без специальной технической подготовки к данным, хранящимся в реляционных базах данных. Показано, что традиционный подход, основанный на ручном написании SQL-запросов, создаёт высокий порог входа и увеличивает зависимость бизнес-пользователей от ИТ-специалистов.

В рамках работы были проанализированы современные подходы к решению задачи Text-to-SQL, включая методы на основе больших языковых моделей, RAG-архитектуры, связывания схемы базы данных и механизмов самокоррекции. На основе проведённого анализа была спроектирована архитектура интеллектуального модуля, включающая сервисы векторизации, индексации документов, поиска релевантного контекста, доступа к базе данных и агентной обработки пользовательского запроса.

В результате была описана программная реализация микросервисной системы, использующей векторное хранилище для работы с метаинформацией схемы данных и агентный цикл для генерации, выполнения и корректировки SQL-запросов. Отдельное внимание уделено обработке ошибок, ограничениям безопасности и локальному выполнению компонентов, что важно для корпоративных сценариев применения.

Поставленная цель работы достигнута: разработано и описано программное решение, позволяющее преобразовывать запросы на естественном языке в SQL-запросы и выполнять интерактивный анализ данных в реляционной базе. Задачи работы выполнены в полном объёме: проведён анализ предметной области и существующих решений, обоснована архитектура, описаны ключевые алгоритмы, приведена программная реализация и рассмотрены результаты применения разработанного подхода.
